% ---------------------------------------------------------------------------
\section{A model of summarising lenses}
\label{sec:model}

\subsection{From selection, function and join to seven stages}

Tominski et al.\ model a lens as a small pipeline embedded in the
visualization pipeline: a selection $\Sel$ picks out the part of the data the
lens affects, a lens function $\lambda$ transforms it, and a join $\Jn$
composes the result with the base visualization~\cite{tominski2014star,
tominski2017survey}. The model is deliberately general, and for the lenses it
was built around---magnifiers, filters, restylers---$\lambda$ is a
transformation of the selected elements and $\Jn$ is compositing into the
lens's interior. Neither needs further structure.

A \emph{summarising} lens is different in both places. Its $\lambda$ does not
transform the selected elements; it replaces them with a much smaller derived
dataset---one value per bin---and then with marks for those values. And its
$\Jn$ is no longer compositing: the marks have to go \emph{somewhere}, and
where they go carries meaning. We therefore refine the two opaque stages into
three each, and keep $\Sel$:
\begin{equation}
  \lambda \;=\; \Mrk \circ \Nrm \circ \Bin,
  \qquad
  \Jn \;=\; \Asc \circ \Anc \circ \Plc .
  \label{eq:refine}
\end{equation}
Binning $\Bin$ decomposes the selected members; normalisation $\Nrm$ turns each
bin into a value against a baseline; marks $\Mrk$ turn values into glyphs.
Placement $\Plc$ decides where on a curve each glyph goes; the anchor $\Anc$
decides what that curve is; association $\Asc$ decides how a glyph is tied back
to the ground it summarises. \Cref{fig:pipeline} shows the seven stages and the
correspondence. The rest of this section defines each stage precisely enough
that the techniques of \cref{sec:descriptive} can be written as paths through
it.

\begin{figure*}[t]
  \centering
  \begin{tikzpicture}[
      node distance=4.2mm,
      stage/.style={draw, rounded corners=2pt, minimum height=9mm, minimum width=20mm,
        align=center, font=\small, fill=white},
      arr/.style={-{Latex[length=2mm]}, thick},
      brace/.style={decorate, decoration={brace, amplitude=5pt, raise=2pt}, thick},
    ]
    \node[stage] (sel) {$\Sel$\\selection};
    \node[stage, right=of sel] (bin) {$\Bin$\\binning};
    \node[stage, right=of bin] (nrm) {$\Nrm$\\normalisation};
    \node[stage, right=of nrm] (plc) {$\Plc$\\placement};
    \node[stage, right=of plc] (anc) {$\Anc$\\anchor};
    \node[stage, right=of anc] (mrk) {$\Mrk$\\marks};
    \node[stage, right=of mrk] (asc) {$\Asc$\\association};
    \draw[arr] (sel) -- (bin);
    \draw[arr] (bin) -- (nrm);
    \draw[arr] (nrm) -- (plc);
    \draw[arr] (plc) -- (anc);
    \draw[arr] (anc) -- (mrk);
    \draw[arr] (mrk) -- (asc);
    \draw[brace] (sel.north west) -- node[above=7pt, font=\small] {$\Sel$} (sel.north east);
    \draw[brace] (bin.north west) -- node[above=7pt, font=\small] {part of $\lambda$} (nrm.north east);
    \draw[brace] (plc.north west) -- node[above=7pt, font=\small] {$\Jn$} (anc.north east);
    \draw[brace] (mrk.north west) -- node[above=7pt, font=\small] {$\lambda$ (marks) \;/\; $\Jn$ (association)} (asc.north east);
    \node[below=2.5mm of nrm, font=\footnotesize, text width=48mm, align=center, text=black!65]
      {what is summarised, and against what};
    \node[below=2.5mm of anc, font=\footnotesize, text width=48mm, align=center, text=black!65]
      {where the summary goes, and on what};
  \end{tikzpicture}
  \caption{The seven stages of a summarising lens, and how they refine
  Tominski et al.'s selection $\Sel$, lens function $\lambda$ and join
  $\Jn$~\cite{tominski2014star}. The execution order interleaves the two
  refinements (marks are sized before they are drawn on the anchor), which is
  why $\lambda=\Mrk\circ\Nrm\circ\Bin$ and $\Jn=\Asc\circ\Anc\circ\Plc$ are
  written as a factorisation rather than as two contiguous blocks.}
  \label{fig:pipeline}
\end{figure*}

\subsection{Selection $\Sel$}

A selection is a region $R$ with an anchor point $c$ from which bearing and
distance are measured. It returns the members $M=\{m_i\}$ inside $R$, each
annotated with its distance $d_i$ and bearing $\theta_i$ from $c$---or, for a
corridor along a polyline $P$, with its chainage $s_i$ (distance along $P$) and
signed offset $o_i$ (distance across, positive to the left of travel). We
distinguish five shapes: a \emph{disc}, an \emph{annulus}, a \emph{sector}, a
\emph{corridor}, and an arbitrary \emph{polygon}. A lasso, an administrative
unit and an isochrone are all polygons: they differ only in where the shape
came from, so a lens that accepts rings accepts all three. A polygon has no
natural centre, so $c$ defaults to its area-weighted centroid and can be
overridden (an isochrone's origin is the obvious override).

For areal data---census units rather than points---a member can be partly
inside $R$. Two standard answers are supported: count a unit if its centroid is
inside, or weight it by the fraction of its area inside. The measure being
aggregated must then declare whether it is \emph{extensive} (counts, which
apportion and sum) or \emph{intensive} (rates, shares and medians, which do
neither and are averaged with a weight). Nothing about a number says which it
is, and summing a column of percentages produces a plausible and wrong
picture, so the implementation requires the declaration.

\subsection{Binning $\Bin$}

Binning partitions $M$ into bins $B_1,\dots,B_K$ and gives each bin a
\emph{preferred position} $p_k\in[0,1)$ on the anchor curve and, optionally, a
\emph{feasible interval} $I_k\subseteq[0,1)$. These two quantities are the only
things placement consumes. Six modes are defined:

\begin{description}[leftmargin=0pt, itemsep=1pt]
  \item[categorical] one bin per category; $p_k$ is a nominal slot, so angle
    carries order and nothing else.
  \item[angular] one bin per bearing sector; $p_k$ is the sector's bearing
    over $360^\circ$, so angle \emph{is} bearing.
  \item[radial] one bin per distance band; there is no bearing, so bins take
    nominal slots and the radial axis is the reading.
  \item[cross] bearing sector $\times$ category.
  \item[chainage] one bin per stretch of a corridor; $p_k = s/L$ on an open
    curve.
  \item[unit] one bin per areal unit; $I_k$ is the arc the unit subtends seen
    from $c$, which is Speckmann and Verbeek's original necklace primitive
    \cite{speckmann2010necklace}.
\end{description}

A categorical bin may also be positioned at the circular mean of its members'
bearings, which is what morphs a category ring into a bearing-faithful one
(\cref{sec:generative}). That position is only meaningful when the members are
concentrated, so every bin also carries the mean resultant length
$R_k\in[0,1]$ of its bearings~\cite{mardia1999directional}; we return to this in
\cref{sec:within}.

\subsection{Normalisation $\Nrm$}

Normalisation maps a bin to a value $v_k$. \emph{Count} and \emph{density}
(count per unit area) are the familiar pair; a radius slider silently
confounds them, which is why normalisation is a stage and not a formatting
option. \emph{Share} is $n_k/n$. The \emph{location quotient}
$\mathrm{LQ}_k=(n_k/n)/(b_k/b)$ and a $z$-score compare the lens against a
baseline profile $b$; \emph{delta} is a signed difference from a reference
(another lens, time or region)---in the terms of visual comparison, an explicit
encoding of the difference rather than juxtaposition or
superposition~\cite{gleicher2011comparison}. The baseline is itself a choice---the
complement of the lens, the study area, or a fixed reference---and changes the
reading materially. One lens defaults to its complement (``what is this place
unusually full of, compared with its surroundings''); a field of lenses must
share one baseline, or every cell becomes average by construction
(\cref{sec:continuum}).

\subsection{Placement $\Plc$}
\label{sec:placement}

Let mark $k$ need half-width $h_k$ on the curve (its own width plus any room
reserved for its label, in units of the curve parameter), and let
$\delta(x,y)$ be the cyclic difference on $[0,1)$. Placement solves
\begin{equation}
  \begin{aligned}
    &\min_{x}\; \sum_{k} w_k\, \delta(x_k,p_k)^2 \\
    &\text{s.t. the arcs } [x_k-h_k,\,x_k+h_k] \text{ are disjoint},\;
    x_k\in I_k, \\
    &\phantom{\text{s.t. }}\text{and the marks keep the cyclic order of the } p_k,
  \end{aligned}
  \label{eq:placement}
\end{equation}
a squared-displacement variant of the necklace-map placement problem of
Speckmann and Verbeek~\cite{speckmann2010necklace,speckmann2015algorithms}. The
last constraint, which we call \emph{order preservation}, corresponds to their
fixed-order case. It also forbids a mark from winding past its neighbours: in the
frame cut open at some point of the cycle, the unwrapped positions increase in
the same order as the unwrapped preferred positions. With order fixed and no
intervals, substituting $y_k=x_k-c_k$, where $c_k$ accumulates the minimum
spacings $h_{k-1}+h_k$, turns non-overlap into monotonicity of $y$. The
optimum is then a weighted isotonic regression, solved exactly in $O(K)$ by
pool-adjacent-violators. The implementation tries each of the $K$ places to cut
the cycle and keeps the cheapest, for $O(K^2)$ overall.

Order preservation is a design constraint, not a free property of the
optimum. We had first assumed, as the implementation's own comments did, that a
crossing solution can always be uncrossed without increasing cost. That holds
when all marks have equal widths and equal weights, and it fails otherwise. We
ran an exact comparison (\texttt{paper/scripts/solver-check.mjs}) on random
instances with $K = 3$--$5$ and collision-prone, clustered preferred positions.
When reordering and winding were allowed, the optimum was cheaper than the
order-preserving one in 8--42\% of instances with unequal widths or unequal
weights, and in none with both equal. The median cost of keeping order was
3--7\textdegree{} of weighted RMS displacement on a closed ring, and at most
31\textdegree. One three-mark example has marks filling 79\% of the ring: a narrow,
heavy mark and a wide one whose preferred positions differ by 2\textdegree, and
a third mark about 30\textdegree{} clockwise of both. Placing the wide mark
counterclockwise of the narrow one, against their bearing order, frees room for
the third mark and nearly halves the cost. For the related max-scale problem, Speckmann and
Verbeek likewise show that ordering by region bearing is not generally optimal
(it is a tight $\frac12$-approximation)~\cite{speckmann2015algorithms}. We keep
the constraint because the relative order of the marks is itself part of what a
bearing-faithful ring encodes. A pair of marks drawn in the wrong order
misstates which of two directions lies clockwise of the other, however small the
displacement saved. Whether readers would accept a small reordering for a
smaller displacement is an empirical question, and could be added to the
study's displacement factor (\cref{sec:evaluation}).

The wrap-around constraint and feasible intervals need no approximation. When the
marks nearly fill the cycle, the wrap-around constraint caps the span of $y$ for
each cut, but the cap never decides the result. The order-preserving optimum
$x^\ast$ has at least one gap wider than its two marks need, because the marks
fill less than the whole curve. At the cut through that gap, $x^\ast$ is feasible
and the cap is slack, so the uncapped isotonic regression for that cut returns
$x^\ast$. Every other cut returns $x^\ast$ or a feasible point that costs more.
(This assumes displacements below half a turn, where cyclic and unwrapped
distances agree.) Intervals $I_k$ become per-mark bounds on $y$, and the same
argument holds. Each cut is then a bounded isotonic regression, solved exactly in
$O(K)$. Monotonicity lets the bounds be tightened to a running maximum from the
left and a running minimum from the right. Pool-adjacent-violators then gives
each pooled block its weighted mean, clamped to the block's feasible range.
Clipping the unconstrained fit to the bounds is \emph{not} enough: for targets
$(0,10,5)$ with $y_2\le 6$ it gives $(0,6,7.5)$, and the optimum is
$(0,6,6)$. A mark wider than its own interval has that interval dropped. If
the remaining intervals cannot all be met, the marks are placed without
intervals, and the solver reports that it did so.

This replaced a heuristic that clamped marks into their intervals and
re-solved. We compared both with an exact reference, Dykstra's algorithm run to
convergence on every cut (\texttt{paper/scripts/interval-check.mjs}). The test
cases were wedge intervals, as angular binning produces, and overlapping arcs,
as areal units produce, for $K = 4$--$16$. With wedges, the heuristic left an
interval in 2--13\% of instances and overlapped a neighbour in 7--13\%. In
37--72\% of instances it returned a valid placement that cost more than the
optimum, by a median of about 5\% and at most 33\%. With overlapping arcs it overlapped in
over half of the instances. The bounded solver matched the reference on every
feasible instance and never overlapped. Without intervals, the solver matched an
exact span-capped reference on all 1,800 closed-ring instances of the order
check.

Two placements are degenerate cases of \cref{eq:placement}: \emph{block}
placement, in which $p_k$ are nominal slots with fixed spacing and the solver
has nothing to do, and \emph{stacked} placement, in which each variable gets
its own concentric curve and competes only with itself, so displacement
vanishes at the cost of radial space and cross-variable comparability.

\subsection{Anchor $\Anc$}

Placement is solved in the curve's parameter, and each mark's room is reserved
as a \emph{fraction} of the curve. Nothing in it assumes a circle. Any curve
of the same length $\ell$ accepts the same solution, so the anchor is a free
axis, and the useful thing to vary along it is curvature. For a ring of radius
$r_0$ ($\ell = 2\pi r_0$) we define the family $\Anc_u$, $u\in[0,1]$: the arc of
length $\ell$ and curvature $(1-u)/r_0$ through a fixed parameter $t_0$,
rotated so that $\Anc_1$ is horizontal. $u=0$ is the closed ring; $u=1$ is a
straight baseline on which bearing becomes a horizontal axis. Changing $u$
never re-runs placement: it is a change of anchor, not a re-solve, so it
composes with every other stage. The seam falls opposite $t_0$; choosing $t_0$
in the widest gap between marks means opening the ring never cuts a mark in
two.

The same axis applies to an open curve. Straightening a corridor onto its own
chainage produces a route profile, which is a \emph{linear cartogram}: chainage
and offset survive exactly, position does not.

The trade the anchor controls is exact and worth stating. A ring keeps
\emph{adjacency}---each mark points at the part of the map it summarises---and
pays for it in comparability, because every mark grows from its own baseline in
its own direction. A straight axis makes lengths directly comparable and gives
labels somewhere to go, and pays in adjacency.

\subsection{Marks $\Mrk$}

A mark turns $v_k$ into a glyph at $\Anc(x_k)$: a radial \emph{bar} (length
$\propto v_k$), a \emph{disc} (area $\propto v_k$, the classic necklace-map
symbol), or a \emph{rose} (the bin's own angular histogram, oriented to true
north so that roses compare with each other and with the compass). Which way a
mark grows is a second, weaker way of buying comparability: along the curve's
outward normal (maximal association), always screen-up (one shared baseline;
on a closed ring the lower marks grow back across the lens), or upright but
never inward (a shared axis with two baselines). A nested glyph such as a rose
carries two readings and must say which owns size; roses default to equal
size, because a rose scaled by value is illegible for exactly the small bins
whose direction is least certain.

\subsection{Association $\Asc$}
\label{sec:association}

On a closed ring around its own selection, adjacency associates a mark with its
ground for free. Two things spend that: \emph{placement}, which moves a mark
off its bearing to avoid overlap, and the \emph{anchor}, which under unrolling
detaches the whole chart from the geography. They are the same loss by
degrees, so one encoding answers both. A \emph{leader} runs from the mark to
the position placement tried to honour---the preferred bearing---at the
members' own mean distance from $c$, drawn in the lens's azimuthal frame. Its
length is therefore the association that was given away.

The leader is triggered by the \emph{gap} between where a mark is and where its
data is, not by the solver's reported displacement. The distinction matters
under block placement: the solver is asked for a nominal slot, reports zero
displacement, and yet every mark is as far from its data as it can be. With
the gap as trigger, morphing a ring from bearing back to category order fades
the leaders in exactly as the marks leave their bearings. A leader is refused
where it would have to guess: a distance band or a nominal slot has no bearing
to point along, and a lens without a metric scale has no distance to point at.

\subsection{Within-unit structure}
\label{sec:within}

Every stage above yields one number per bin, and that number hides whether the
bin's members are spread evenly or piled in one corner---which is what makes an
aggregate fragile to how its unit was drawn. Hexagon-based techniques have
begun to encode this: Honeycomb Plots cut each tile along the regression plane
of its within-tile density, a ``diamond cut''~\cite{trautner2022honeycomb}, and
HexTiles derive a per-tile confidence from weighted within-tile variance, argued
explicitly as mitigation of the modifiable areal unit
problem~\cite{kawakami2024hextiles}.

A lens gets this more cheaply than a tile, for a structural reason. A hexagon
has no privileged origin or orientation, so the within-tile distribution must be
fitted. A lens is already a polar coordinate system centred on a point the
analyst chose, so its within-unit distribution decomposes, with no fitting,
into bearing and distance. For each bin we compute the circular mean
$\bar\theta_k$ and mean resultant length $R_k$ of member bearings, the circular
standard deviation $\sqrt{-2\ln R_k}$~\cite{mardia1999directional}, and the
mean, coefficient of variation and edge share of member distances; on a
corridor, the signed mean offset (\emph{bias}) and the mean offset sign
(\emph{sidedness}), which disagree precisely when provision is balanced by
count but not by distance. These are drawn as a \emph{spread} arc (circular SD),
a \emph{gradient} arrow (lens-level $\bar\theta$, length $\propto R$), the
members drawn back through the aggregate (\emph{inclusions}, after Honeycomb's
``amber inclusions''), or the rose mark.

\subsection{The radius control as a scale diagnostic}
\label{sec:elasticity}

The unit of a lens is modifiable by a control the analyst is already holding,
so the lens can report how much its reading depends on that choice. Let $N(r)$
be the number of members within distance $r$ of $c$. We define the
\emph{elasticity} of the count with respect to the radius,
\begin{equation}
  E(r) \;=\; \frac{d\ln N}{d\ln r} \;=\; \frac{r\,N'(r)}{N(r)}
  \;=\; 2\,\frac{\lambda_{\mathrm{rim}}(r)}{\bar\lambda(r)},
  \label{eq:elasticity}
\end{equation}
where $\lambda_{\mathrm{rim}}=N'(r)/2\pi r$ is the intensity on the rim and
$\bar\lambda=N/\pi r^2$ the mean intensity inside. The last form follows
directly and is the useful reading: $E=2$ when the rim is as dense as the
interior (uniform density), $E\to 0$ when everything is already inside, and
$E\gg2$ when a cluster sits just beyond the rim and the reading is about to
jump. The implementation uses a backward difference over a band of relative
width $b$,
\begin{equation}
  \hat E(r) \;=\; \frac{N(r)-N\big((1-b)r\big)}{b\,N(r)},
  \label{eq:elasticity-hat}
\end{equation}
which tends to $E$ as $b\to0$ and equals $2-b$ (not $2$) under uniform density;
with the default $b=0.1$ the uniform reference is $1.9$. Samples with fewer than
a floor of members (30 by default) are flagged unreliable, because
\cref{eq:elasticity-hat} is a ratio of small counts.

$E(r)$ does not depend on the radius currently set, so it is computed once per
centre and drawn on the radius slider itself: the analyst sees where the cliffs
are before dragging onto one (\cref{fig:elasticity}). This is an instance of
embedding a visualization in the control that it informs, in the manner of
scented widgets~\cite{willett2007scented}.

We do not claim the statistic. Counting within growing distances of a point is
the idea behind Ripley's $K$ function, for which $K(t)=\pi t^2$ under a Poisson
process~\cite{ripley1977modelling}, and behind second-order neighbourhood
analysis, which takes that view from each individual
point~\cite{getis1987second}; under complete spatial randomness $N\propto r^2$
and so $E=2$. The same log--log slope around a freely chosen centre is the
mass--radius (fractal) dimension of radial analysis in urban
morphology~\cite{yamu2016fractal}. Three things distinguish the lens's use of it:
the centre is an arbitrary location rather than an event, so this is a
point-to-event rather than an event-to-event count; the uniform reference
depends on the lens shape (it is 2 for a disc or sector, but not for an annulus
with fixed inner radius, nor 2 for a corridor whose width is the control); and
it is shown as a property of the control rather than of the result.
The library therefore reports $\hat E$ with two references in place of the bare
value 2. The first is the estimator's own value under uniform density for the
selection's shape, $\big(A(r)-A((1-b)r)\big)/\big(b\,A(r)\big)$ for a selection
of area $A(r)$. That is $2-b$ for a disc or sector and larger for an annulus
with a fixed inner radius, and the ratio of $\hat E$ to it is 1 under uniform
density for any shape. The second is a pointwise Monte-Carlo envelope under
complete spatial randomness, conditioned on the number of members and in the
spirit of the simulation envelopes used with Ripley's
$K$~\cite{ripley1977modelling}. It widens where counts are small, which makes
it the principled version of the 30-member floor.
\doubt{The equivalences are derived here. Ripley and Yamu et al.\ were read in
full; Getis \& Franklin at abstract level only. Yamu et al.\ is mainly a
planning-theory paper that uses radial analysis; a methods paper by Frankhauser
would be the more canonical citation for mass--radius dimension (not yet
identified and checked).}

\begin{figure}[t]
  \centering
  \begin{tikzpicture}
    \begin{axis}[
        width=\columnwidth, height=45mm,
        xlabel={radius $r$ (m)}, ylabel={$\hat E(r)$},
        xmin=0, xmax=3000, ymin=0, ymax=8,
        ytick={0,1.9,4,6,8}, yticklabels={0,1.9,4,6,8},
        grid=major, grid style={black!8},
        tick label style={font=\footnotesize}, label style={font=\footnotesize},
        axis lines*=left,
      ]
      \addplot[black!25, fill=black!6, draw=none] coordinates {(0,0) (0,8) (475,8) (475,0)} \closedcycle;
      \addplot[name path=lo, draw=none, unbounded coords=jump]
        table[x=r, y expr={\thisrow{reliable}>0 ? \thisrow{low} : nan}, col sep=comma]
        {figures/elasticity.csv};
      \addplot[name path=hi, draw=none, unbounded coords=jump]
        table[x=r, y expr={\thisrow{reliable}>0 ? \thisrow{high} : nan}, col sep=comma]
        {figures/elasticity.csv};
      \addplot[RoyalBlue!15] fill between[of=lo and hi, soft clip={domain=475:3000}];
      \addplot[thick, color=RoyalBlue, unbounded coords=jump]
        table[x=r, y expr={\thisrow{reliable}>0 ? \thisrow{elasticity} : nan}, col sep=comma]
        {figures/elasticity.csv};
      \addplot[dashed, black!60, domain=0:3000, samples=2] {1.9};
      \node[font=\scriptsize, anchor=north west, align=left, text=black!55] at (axis cs:30,7.8) {too few\\members};
      \node[font=\scriptsize, anchor=south west, text=black!60] at (axis cs:2100,2.2) {uniform: $2-b$};
    \end{axis}
  \end{tikzpicture}
  \caption{The elasticity profile $\hat E(r)$ (\cref{eq:elasticity-hat}, $b=0.1$)
  around the demonstration centre in Yogyakarta, computed from the bundled
  OpenStreetMap extract by the library and plotted from its output. The band is
  a pointwise 95\% Monte-Carlo envelope under complete spatial randomness: 999
  patterns with the same 1{,}071 members within 3\,km, scattered uniformly over
  the disc. The shaded span has fewer than 30 members and is not drawn. The
  early cliff ($\hat E$ up to 7.3 at 475--525\,m, against an envelope reaching
  3.6) says a band of places at that distance is far denser than anything
  inside it, so a radius chosen there is fragile. The rise at 925--1000\,m
  stays inside the envelope, which is what a uniform pattern of this size
  produces by chance. From 1.5 to 2.7\,km the curve lies below the envelope
  throughout: widening the lens there adds ground faster than it adds places.
  The envelope is pointwise, so a brief excursion at one of the 102 drawn radii is
  weak evidence on its own.}
  \label{fig:elasticity}
\end{figure}
